% PREAMBULO DE LA VERSION FINAL 2 (copia del de capitulos_final1/, 2026-09-23).
% Cada version final lleva su preambulo, su bibliografia y sus figuras rehechas dentro de
% capitulos_finalN/, para que corregirla no altere la version anterior.
%
% =====================================================================
%  preambulo.tex  -- Paquetes y configuracion de la tesis EduFEM
%  Cargado por las DOS versiones del documento. Clase: report (12pt, a4paper).
%
%    main.tex     pdflatex + biber    Latin Modern (serif)   -> main.pdf
%    main_v2.tex  xelatex  + biber    Arial 12 pt            -> main_v2.pdf
%
%  Mismos capitulos, misma bibliografia, mismo formato APA: lo unico que cambia
%  es la fuente de TEXTO. Editar un capitulo actualiza las dos versiones.
% =====================================================================

% --- Codificacion y fuente de TEXTO ---
% main_v2.tex define \tesisFuenteArial ANTES de este \input; ese es el unico
% interruptor de todo el mecanismo.
%
% Las FORMULAS no se tocan en ninguna de las dos versiones: amsmath las sigue
% componiendo con las fuentes matematicas de LaTeX (Latin Modern Math). Es lo que
% corresponde en una tesis de ingenieria y ademas no hay alternativa real: Arial
% no trae alfabeto matematico, ni italica matematica, ni los simbolos de extension
% (integral, sumatoria, \nabla, los delimitadores que crecen con la matriz).
% Forzarla romperia las ecuaciones del marco teorico y del Anexo F.
\ifdefined\tesisFuenteArial
  % --- XeLaTeX: la Arial de verdad, la del sistema (C:\Windows\Fonts\arial.ttf) ---
  % lmodern va PRIMERO y a proposito: declara las familias MATEMATICAS en Latin
  % Modern, las mismas de main.pdf. Recien despues fontspec redefine las familias
  % de TEXTO (rm/sf/tt), asi que las formulas salen identicas a las de la v1.
  \usepackage{lmodern}
  \usepackage{fontspec}
  % Ligatures=TeX (-- , --- y las comillas de TeX) va SOLO en las familias de
  % texto, nunca en la monoespaciada: en un listado de codigo convierte el "--"
  % de Python en una raya y las comillas rectas en curvas.
  \setmainfont{Arial}[Ligatures=TeX]
  \setsansfont{Arial}[Ligatures=TeX]
  % La monoespaciada se queda en Latin Modern Mono, la MISMA de main.pdf, y no
  % pasa a una mono de Windows. No es cosmetica: medido sobre el Anexo C, con
  % Consolas el codigo sale con U+2010 HYPHEN en vez del guion ASCII (y con
  % Ligatures=TeX, ademas con U+201D en vez de la comilla recta), asi que lo que
  % el tribunal copia del PDF ya no es Python valido. Con Latin Modern Mono el
  % texto extraido no tiene un solo caracter espurio, igual que en la v1.
  \setmonofont{Latin Modern Mono}
\else
  % --- pdflatex: la version original ---
  \usepackage[utf8]{inputenc}      % fuente .tex en UTF-8
  \usepackage[T1]{fontenc}
  \usepackage{lmodern}
\fi
% es-noindentfirst se retiro al adoptar APA: la norma sangra la primera linea de
% TODO parrafo, incluido el primero que sigue a un titulo.
\usepackage[spanish,es-tabla]{babel}  % "Tabla" (no "Cuadro")

% --- Parche: babel-spanish rompe los \caption con formulas bajo XeLaTeX ---
% Sintoma exacto: "! Argument of \es@a has an extra }." al procesar un
% \caption{... $D_{\max}$.}, que aborta la compilacion de main_v2.
% Causa: babel-spanish redefine \max, \min, \lim e \inf para que impriman "max",
% "min", "lim" e "inf" ACENTUADOS, que es lo correcto en espanol. Para poner ese
% acento elige entre dos implementaciones segun el motor, y la que reserva para
% XeLaTeX (\es@op@ac@TU, activada porque ahi existe \Umathchar) no sobrevive a la
% expansion que hace \caption al escribir el rotulo en el archivo .lof.
% La rama base -la que usa pdflatex, o sea la que ya produce main.pdf- si.
% Verificado: con este \let, el texto extraido del PDF es identico al de
% pdflatex ("max" "min" "lim" "inf" con tilde); sin el, no compila. La
% alternativa \unaccentedoperators tambien evita el error pero deja los
% operadores SIN acento, apartandose de main.pdf.
\ifdefined\XeTeXversion
  \makeatletter
  \@ifundefined{es@op@ac@base}{}{\let\es@op@ac\es@op@ac@base}
  \makeatother
\fi
\usepackage{csquotes}            % requerido por biblatex

% babel-spanish redefine \% para anteponer un espacio fino, y esa version rompe en
% modo matematico ("Incompatible glue units"). Se restablece un \% simple, valido en
% texto y en formulas (los capitulos ya escriben el espacio fino con \, donde hace falta).
\AtBeginDocument{\DeclareRobustCommand{\%}{\char37\relax}}

% Guion bajo que admite corte de linea justo despues de el. Permite que los
% identificadores de codigo largos (p.ej. generate_structured_quad_mesh,
% boundary_conditions.csv) se ajusten dentro de las celdas p{} estrechas de las
% tablas de anexos sin desbordar el margen. Es puramente aditivo: solo agrega
% oportunidades de corte, nunca fuerza un salto.
\AtBeginDocument{\DeclareRobustCommand{\_}{\textunderscore\allowbreak}}

% =====================================================================
%  PRESENTACION DEL DOCUMENTO: formato APA 7.a edicion
% ---------------------------------------------------------------------
%  Las CITAS y la lista de REFERENCIAS no son APA: siguen Vancouver
%  (numericas, por orden de aparicion). Ver el bloque de biblatex mas
%  abajo y tesis/normas/guia_vancouver.pdf. Lo que sigue es solo la
%  presentacion: margenes, fuente, interlineado, sangria, titulos,
%  rotulos de tablas y figuras y numeracion de pagina.
%  Que se aplico y que no, y por que: docs/notas/2026-09-10_apa-presentacion.md
% =====================================================================

% --- Geometria, interlineado y parrafos (APA 7) ---
% APA pide 2,54 cm (1 pulgada) en los cuatro lados. Se conserva A4 en vez de
% carta: es el papel de la region y el de la caratula ya impresa.
% headheight: fancyhdr necesita 14,5 pt para el numero de pagina a 12 pt; con el
% valor por defecto (12 pt) avisaba en cada hoja.
\usepackage[a4paper,margin=2.54cm,headheight=14.5pt]{geometry}
\usepackage{setspace}
% APA: interlineado doble en todo el documento. Medido sobre esta tesis: a doble
% son 143 hojas y a 1,5 son 120 (el original, antes de APA, eran 131). El salto
% es chico porque el margen de 2,54 cm ensancha la caja de texto y compensa.
% Para pasar a 1,5, cambiar esta linea por \onehalfspacing.
\doublespacing
\setlength{\parskip}{0pt}        % APA: sin espacio extra antes ni despues del parrafo
\setlength{\parindent}{1.27cm}   % APA: sangria de primera linea de 1/2 pulgada

% Excepciones al doble espacio que la propia APA admite: el cuerpo de las tablas
% y el texto dentro de figuras "puede utilizar interlineado sencillo, 1,5 o doble,
% el que sea mas efectivo". Se agregan los listados de codigo por el mismo motivo:
% a doble espacio un fragmento de Python se vuelve ilegible.
% APA: "Align the text to the left margin. Leave the right margin uneven, or
% ragged. Do not use full justification." Y sin guiones de corte automaticos.
% Para volver a texto justificado, comentar estas tres lineas.
\usepackage[document]{ragged2e}
% ragged2e pone la sangria de primera linea en 0 al activar \RaggedRight; APA la
% quiere en 1,27 cm, asi que se restituye por su propio parametro. OJO: este
% parametro tambien entra en las celdas p{} de las tablas, donde no va; se lo
% neutraliza mas abajo, en el parche de \@arrayparboxrestore.
\setlength{\RaggedRightParindent}{1.27cm}
% Bandera derecha LIBRE. ragged2e limita por defecto el estiramiento del margen
% derecho a 2 em, y con los guiones prohibidos eso deja a TeX sin salida: no
% puede cortar la palabra ni correrla al renglon siguiente, asi que la mete en
% el margen. Eran 254 lineas desbordadas, la peor por 1,7 cm. Con "plus 1fil"
% el renglon simplemente termina antes, que es la bandera derecha de verdad y
% lo que APA pide ("leave the right margin uneven, or ragged").
\setlength{\RaggedRightRightskip}{0pt plus 1fil}
\hyphenpenalty=10000
\exhyphenpenalty=10000

\usepackage{etoolbox}
\AtBeginEnvironment{tabular}{\singlespacing}
\AtBeginEnvironment{tabularx}{\singlespacing}
\AtBeginEnvironment{longtable}{\singlespacing}
\AtBeginEnvironment{lstlisting}{\singlespacing}

% --- Las celdas p{} NO heredan la sangria ni la prohibicion de guiones ---
% La opcion [document] de ragged2e activa "raggedrightboxes", que reemplaza el
% \parindent\z@ de \@arrayparboxrestore por un \RaggedRight. Como \RaggedRight
% aplica \RaggedRightParindent, TODA celda p{} quedaba con 1,27 cm de sangria en
% su primera linea: el encabezado y la primera linea de cada celda se corrian a
% la derecha y las lineas siguientes quedaban al ras, que es exactamente el
% desalineado que se veia en la Nomenclatura y en la tabla de variables.
% Ademas, sin guiones una columna angosta no puede partir "caracteristico" ni
% "discretizacion" y se desborda. Adentro de una tabla la norma APA no manda:
% permite explicitamente que el cuerpo se componga como resulte mas efectivo.
\makeatletter
\apptocmd{\@arrayparboxrestore}{%
  \parindent\z@
  \hyphenpenalty=50 \exhyphenpenalty=50 % los valores por defecto de LaTeX
}{}{\PackageError{tesis}{No se pudo parchear \string\@arrayparboxrestore}{}}

% --- Los indices (contenido, figuras, tablas) van justificados ---
% Sus renglones terminan en un relleno de puntos y el numero de pagina; con la
% bandera derecha libre ese relleno pelea con el \rightskip y el titulo de la
% figura se mete en el margen. No son texto corrido, asi que la alineacion de
% APA no les aplica.
% Y ADEMAS admiten guion de corte, por el mismo motivo que las celdas p{} de
% mas arriba: el titulo de una figura o de una seccion no se puede reformular
% para que entre, y aca no hay bandera derecha que absorba el sobrante -el
% relleno de puntos ocupa hasta el numero de pagina-, asi que sin guiones la
% unica salida que le queda a TeX es meter la palabra en el margen. Medido: con
% Arial se salian 8 pt (2,8 mm) la entrada del capitulo 2 y una de la lista de
% figuras. Los valores son los que LaTeX trae por defecto.
\let\tesis@starttoc\@starttoc
\renewcommand{\@starttoc}[1]{\begingroup\justifying
  \hyphenpenalty=50 \exhyphenpenalty=50
  \tesis@starttoc{#1}\endgroup}
\makeatother

% --- Matematica ---
\usepackage{amsmath,amssymb,amsfonts}
\usepackage{bm}                  % \bm para negrita matematica (matrices)

% Operador de ensamblaje del MEF: la "A grande" estandar (no el wedge logico).
% Con \DeclareMathOperator* el subindice (el indice de elemento) se ubica debajo.
\DeclareMathOperator*{\Ensamblaje}{\text{\Large\textsf{A}}}

% --- Figuras y tablas ---
\usepackage{graphicx}
\graphicspath{{capitulos_final2/figuras/}{figuras/}}   % final2: las figuras rehechas tienen prioridad
\usepackage{booktabs}
\usepackage{array}
\usepackage{longtable}
\usepackage{multirow}
\usepackage{float}
\usepackage{caption}
% --- Rotulos de tablas y figuras en formato APA 7 ---
% APA: numero en negrita en su propio renglon, titulo en cursiva debajo, ambos
% alineados a la izquierda; la nota va debajo del contenido (comando \notaapa).
\captionsetup{labelfont=bf,        % "Figura 1.1" en negrita
              textfont=it,         % titulo en cursiva
              labelsep=newline,    % titulo en el renglon siguiente
              singlelinecheck=false,
              justification=raggedright,
              font=normalsize}

% APA coloca numero y titulo ARRIBA de la figura, no debajo. plaintop lo resuelve
% sin tocar los 29 entornos figure de los capitulos, que traen el \caption al final.
\floatstyle{plaintop}
\restylefloat{figure}

% Nota al pie de una tabla o figura, en el formato de APA: "Nota." en cursiva.
% Uso: dentro del float, despues del contenido.  \notaapa{Elaboracion propia.}
\newcommand{\notaapa}[1]{%
  \par\vspace{0.4em}%
  {\raggedright\footnotesize\textit{Nota.} #1\par}}
\usepackage{pdfpages}            % embeber PDF externo (Anexo de validacion SAP2000)

% --- Color, enlaces y tipografia fina ---
\usepackage{xcolor}
\definecolor{tesisAzul}{HTML}{0D6EFD}
\usepackage{microtype}
\usepackage[hidelinks]{hyperref}
\hypersetup{
  pdftitle={\tituloTesis},
  pdfauthor={\autorTesis},
  pdfsubject={Tesis de licenciatura en Ingeniería Civil, Universidad Autónoma Tomás Frías},
  pdfkeywords={método de elementos finitos, software educativo, trazabilidad, verificación y validación, memoria de cálculo},
  bookmarksnumbered=true,
}
% \autoref{anexo:...} debe imprimir "Anexo X" (el rotulo de los capitulos del apendice)
% tambien ANTES de \appendix (la Introduccion los enumera). babel-spanish restablece
% los nombres de hyperref en cada seleccion de idioma, por eso se agrega a \extrasspanish.
\addto\extrasspanish{\renewcommand{\appendixautorefname}{Anexo}}

% --- Listados de codigo fuente (Anexo C) ---
% El paquete listings con un mapeo `literate` para los simbolos unicode que
% aparecen en los comentarios del codigo (xi, eta, derivadas parciales, letras
% griegas), de modo que los fragmentos reales del proyecto se reproduzcan fieles.
\usepackage{listings}
\definecolor{lstKeyword}{HTML}{0D6EFD}
\definecolor{lstComment}{HTML}{6C757D}
\definecolor{lstString}{HTML}{198754}
\definecolor{lstBg}{HTML}{F6F8FA}
\lstdefinestyle{pyedufem}{
  language=Python,
  basicstyle=\ttfamily\footnotesize,
  keywordstyle=\color{lstKeyword}\bfseries,
  commentstyle=\color{lstComment}\itshape,
  stringstyle=\color{lstString},
  backgroundcolor=\color{lstBg},
  numbers=left, numberstyle=\tiny\color{lstComment}, numbersep=8pt,
  % Los numeros de linea cuelgan a la izquierda del marco: sin este margen
  % invadian el margen APA de 2,54 cm.
  xleftmargin=1.6em, xrightmargin=0.2em,
  frame=single, framesep=4pt, rulecolor=\color{lstComment},
  breaklines=true, breakatwhitespace=false, showstringspaces=false, upquote=true,
  tabsize=4, columns=fullflexible, keepspaces=true,
  extendedchars=true, inputencoding=utf8,
  literate=%
    {ξ}{{$\xi$}}1 {η}{{$\eta$}}1 {∂}{{$\partial$}}1 {∈}{{$\in$}}1
    {ε}{{$\varepsilon$}}1 {γ}{{$\gamma$}}1 {σ}{{$\sigma$}}1 {τ}{{$\tau$}}1
    {ν}{{$\nu$}}1 {λ}{{$\lambda$}}1 {μ}{{$\mu$}}1 {θ}{{$\theta$}}1 {π}{{$\pi$}}1
    {·}{{$\cdot$}}1 {×}{{$\times$}}1 {→}{{$\rightarrow$}}1 {∇}{{$\nabla$}}1
    {²}{{\textsuperscript{2}}}1 {³}{{\textsuperscript{3}}}1
    {≤}{{$\leq$}}1 {≥}{{$\geq$}}1 {≈}{{$\approx$}}1 {±}{{$\pm$}}1 {√}{{$\surd$}}1
    {—}{{---}}1 {–}{{--}}1 {°}{{$^{\circ}$}}1
    {á}{{\'a}}1 {é}{{\'e}}1 {í}{{\'i}}1 {ó}{{\'o}}1 {ú}{{\'u}}1
    {Á}{{\'A}}1 {É}{{\'E}}1 {Í}{{\'I}}1 {Ó}{{\'O}}1 {Ú}{{\'U}}1
    {ñ}{{\~n}}1 {Ñ}{{\~N}}1 {ü}{{\"u}}1 {¿}{{?`}}1 {¡}{{!`}}1
}
\lstset{style=pyedufem}

% --- Encabezado: numeracion de pagina en formato APA 7 ---
% APA (trabajo de estudiante): el encabezado lleva UNICAMENTE el numero de
% pagina, arriba a la derecha, en todas las hojas y en la misma fuente y cuerpo
% que el texto. Sin titulillo (running head, solo para trabajos profesionales),
% sin nombre de capitulo y sin linea de encabezado.
\usepackage{fancyhdr}
\pagestyle{fancy}
\fancyhf{}
\fancyhead[R]{\thepage}
\renewcommand{\headrulewidth}{0pt}
\fancypagestyle{plain}{% las hojas de apertura de capitulo usan "plain" en report
  \fancyhf{}
  \fancyhead[R]{\thepage}
  \renewcommand{\headrulewidth}{0pt}
}

% --- Jerarquia de titulos en formato APA 7 (conservando la numeracion) ---
% APA define cinco niveles, todos en negrita y con cada palabra en mayuscula.
% APA ademas pide no numerar los titulos; aca SI se numeran (decision del autor):
% la tesis tiene ~50 referencias cruzadas del tipo \autoref{cap:...} y el tribunal
% espera capitulos numerados. Lo que se adopta es la tipografia de cada nivel.
%
%   Nivel 1 (\chapter)       centrado, negrita
%   Nivel 2 (\section)       izquierda, negrita
%   Nivel 3 (\subsection)    izquierda, negrita cursiva
%   Nivel 4 (\subsubsection) sangrado, negrita, punto final, texto en la misma linea
%   Nivel 5 (\paragraph)     sangrado, negrita cursiva, punto final, texto seguido
%
% El cuerpo es el mismo del texto (APA: "utiliza el mismo tamano de fuente en
% todo el documento, incluso en los titulos"), por eso van en \normalsize.
\usepackage{titlesec}

\titleformat{\chapter}[display]
  {\normalfont\normalsize\bfseries\centering}
  {\chaptertitlename\ \thechapter}{0.6\baselineskip}{}
\titlespacing*{\chapter}{0pt}{0pt}{\baselineskip}

% [block] y no la forma "hang" por defecto: con hang, un titulo largo -"1.8
% Fenomenos numericos en elementos de bajo orden: bloqueo y modos espurios"- no
% envuelve y se mete 0,9 cm en el margen derecho. En bloque, parte a dos
% renglones. El \RaggedRight explicito evita que el titulo salga justificado.
\titleformat{\section}[block]
  {\normalfont\normalsize\bfseries\RaggedRight}{\thesection}{1em}{}
\titlespacing*{\section}{0pt}{\baselineskip}{0pt}

\titleformat{\subsection}
  {\normalfont\normalsize\bfseries\itshape}{\thesubsection}{1em}{}
\titlespacing*{\subsection}{0pt}{\baselineskip}{0pt}

\titleformat{\subsubsection}[runin]
  {\normalfont\normalsize\bfseries}{\thesubsubsection}{1em}{}[.]
\titlespacing*{\subsubsection}{\parindent}{\baselineskip}{0.5em}

\titleformat{\paragraph}[runin]
  {\normalfont\normalsize\bfseries\itshape}{}{0pt}{}[.]
\titlespacing*{\paragraph}{\parindent}{\baselineskip}{0.5em}

% --- Sangria francesa de la lista de referencias (APA: 1,27 cm) ---
% Aplica a la lista de biblatex, que es Vancouver en su contenido pero se
% presenta con la sangria que pide APA, como el resto del documento.
\AtBeginDocument{\setlength{\bibhang}{1.27cm}}

% --- Bibliografia estilo VANCOUVER (biblatex + biber) ---
% Numeracion por orden de aparicion (sorting=none), citas numericas compactas.
% numeric-comp + sorting=none + el formato de nombres de abajo reproduce Vancouver/NLM.
\usepackage[backend=biber,    % obligatorio: biber, NO bibtex
            style=numeric-comp,  % numerico con rangos comprimidos: [2-5]
            sorting=none,        % numeracion por orden de aparicion (Vancouver)
            sortcites=true,      % ordena numeros dentro de una cita: [3,1] -> [1,3]
            giveninits=true,     % nombres como iniciales
            terseinits=true,     % iniciales pegadas y sin puntos: "Bathe KJ", "Ryu HR".
                                 % Fija bibinitperiod, bibinitdelim Y bibinithyphendelim;
                                 % este ultimo (por defecto ".-") es el que hacia salir
                                 % "Bathe K.-J." en los nombres compuestos.
            uniquename=false,
            uniquelist=false,
            maxbibnames=6,       % hasta 6 autores en la lista...
            minbibnames=6,       % ...y si hay mas, corta con "et al." (Vancouver/NLM)
            doi=true,            % mostrar DOI en articulos
            isbn=false,
            url=false]{biblatex}

% Formato de nombres estilo Vancouver/NLM: "Salari K" (apellido + iniciales,
% sin coma, sin puntos ni espacios entre iniciales) y sin "y"/"and" final.
% (No hace falta redefinir "sortname": con sorting=none nunca se usa.)
\DeclareNameAlias{default}{family-given}
\renewcommand*{\revsdnamepunct}{}        % sin coma entre apellido e iniciales
% (bibinitperiod y bibinitdelim los deja vacios la opcion terseinits de arriba,
%  junto con bibinithyphendelim, que estos dos \renewcommand no alcanzaban.)
% El ultimo autor se separa con "y" (p.ej. "Salari K y Knupp P"): es la convencion
% de babel-spanish y la natural en un documento en espanol. El NLM ingles usa coma;
% para forzarla habria que sobrescribir el delimitador de babel, lo que es fragil.

\addbibresource{capitulos_final2/referencias.bib}   % final2: copia propia de la bibliografia

% Citas multiples con localizador (\autocites[p.~a]{k1}[p.~b]{k2}): el delimitador
% por defecto (coma) era ambiguo, "[20, pp. 21, 26, 19, pp. 219-222]". Se redeclara
% SOLO el comando multiple \parencites (al que \autocites se mapea con autocite=inline)
% con punto y coma, "[20, pp. 21, 26; 19, pp. 219-222]", y se deja \multicitedelim
% en coma para que \autocite{k1,k2} siga imprimiendo "[4, 14]" como pide Vancouver.
\DeclareMultiCiteCommand{\parencites}[\mkbibbrackets]{\parencite}{\addsemicolon\space}
\ExecuteBibliographyOptions{autocite=inline}

% --- Comando para figuras todavia inexistentes (placeholder compilable) ---
% Uso: dentro de un entorno figure, \figpend{descripcion de la figura}.
% Cuando exista la imagen real, reemplazar por \includegraphics{...}.
\newcommand{\figpend}[1]{%
  \begin{center}
  \setlength{\fboxsep}{10pt}%
  \fcolorbox{tesisAzul}{white}{%
    \parbox[c][3.4cm][c]{0.82\linewidth}{\centering
      \textbf{[Figura pendiente]}\\[4pt]\small\itshape #1}}%
  \end{center}}

% Etiqueta de "pendiente" para texto que falta completar.
\newcommand{\pendiente}[1]{\textcolor{red}{\textbf{[PENDIENTE: #1]}}}

% final1: anchos de los indices (H-219, H-339, H-257)
% El numero de pagina se pegaba al texto en las entradas largas y el numero de figura de dos
% digitos («B.10.») al titulo: se ensancha la caja del numero y el margen derecho de las lineas.
\makeatletter
\renewcommand{\@pnumwidth}{2.2em}
\renewcommand{\@tocrmarg}{3.2em}
\renewcommand*\l@figure{\@dottedtocline{1}{1.5em}{3.0em}}
\renewcommand*\l@table{\@dottedtocline{1}{1.5em}{3.0em}}
\makeatother
